ソーシャルメディアの登場により、インターネット上では多様な投稿が大量に発信されています。
健全なオンライン空間を実現するためには、これらの投稿の背後に込められた「意図」を正確に分析することが重要となります。

私たちは、意図解析に関する国際コンペティション「SemEval 2020」にて、複数のタスクに参加しました。
例えば、「ニュースの記事タイトルを面白おかしく改変したポスト」について、その面白さの度合いを言語モデルで評価するタスクや、ミーム画像に込められた感情を推定するタスク、などです。

私たちはまず、最先端の言語モデルに対して、各タスクの性質に合わせた特化アーキテクチャを追加しました。
例えば前述のユーモア評価タスクの場合は、改変前後の記事タイトルを入力し、新しく追加したcross-attention層によって、両者の差分や関係性を比較・分析させました。
さらに、単一の言語モデルでは捉えきれない特徴を補完するために、複数の言語モデルを独自のレシピに基づいてアンサンブルしました。
その結果、高精度を達成し、\textbf{複数のタスクにおいて1位}を獲得しました。

また、意図を正確に捉えるためには、文中の単語や句がどのような関係を持つのかという、文の意味構造を解析することも重要です。
私たちは、この構造解析、すなわちパージングを題材とする「CoNLL 2020 Shared Task」にも参加し、こちらでも1位を獲得しました。

以上の取り組みを通じて、最先端の言語モデルを実課題に適用する際には、タスクの性質を踏まえたモデル設計や学習方法の工夫が重要であることを学びました。
なかでも、複数モデルの強みを組み合わせるアンサンブル学習の有効性を強く実感し、後の研究テーマへとつながりました。
